\begin{theorem}[The full automorphism group]\hypertarget{tag-01QF}{}\label{fam:src:B.4.1}\label{fam:st:B.4.1}\hypertarget{tag-02YQ}{}\label{night:P5.5}\label{stmt:P5.5}
\StatementRecord{night:P5.5}{proved}{02YQ}{the definitions and constructions stated in the proof}

For every admissible twist of the $(3,4)$ construction,
\[
 \operatorname{Aut}(X_c)=\mathbb C/\mathbb Z\simeq\mathbb C^*.
\]
The parameter $a\bmod\mathbb Z$ acts by $\zeta\mapsto\zeta+ae$,
where $e=\Pi_c\widehat\delta$. The exponential group isomorphism
is $a\bmod\mathbb Z\mapsto e^{2\pi ia}$.
\end{theorem}
\begin{proof}
\begin{enumerate}[label=\arabic*.,leftmargin=*]
\item Proposition~\ref{fam:src:A.1.1} identifies the complete
anticanonical pencil with $f$ and proves that every automorphism
induces the identity on its base.

\item Apply the affine-lift and centralizer calculation of
\cref{fam:src:A.2.1} with $c'=c$. The period equation
\eqref{wb:pshift:eq:period-morphism} gives the integral linear part
$P=aI+bE_{41}$, $a=\pm1$, and the complex linear part $H=aI$;
its first column gives $b=0$.

\item The finite-root equation
\eqref{wb:pshift:eq:finite-norm-obstruction}, with $n=0$,
gives $N_1\ell=-(1-a)v_1$. For $a=-1$, applying the invariant
covector $\gamma$ gives $3\gamma(\ell)=-2l_1$, where
$l_1=\gamma(v_1)$. Freeness gives $3\nmid l_1$, so this
integer equality is impossible. Thus $a=1$, and every
automorphism is fibrewise translation.

\item Let $j$ be the
inclusion of the cusp-punctured orbifold and set
$\mathcal H=j_*\Lambda$.  Its cusp stalk is
$\ker(M_0-I)=\mathbb Z\widehat w\oplus\mathbb Z\widehat\delta$.
Let $\mathcal A$ be the sheaf of fibrewise translations extending to
the local compactification.  There is an exact sequence
\begin{equation}\hypertarget{tag-01OH}{}\label{eq:P5exponential}
0\longrightarrow\mathcal H\longrightarrow E\longrightarrow\mathcal A
\longrightarrow0.
\end{equation}
For the local check at the cusp, use the toroidal covering whose
component lattice is $\Gamma=\Lambda/\Lambda_{\rm tor}$.  A translation
acts identically on the nearby fibre lattice and on $\Gamma$, and
therefore lifts.  Its meridian remains bounding a disc through a smooth
point of $W$, which verifies the same assertion for the full local
cover.  On the punctured torus the lift has form $x\mapsto a(q)x$.
It sends the component $E_0$ to $E_v$; the index-one map
$B_0:\Gamma\to\mathbb Z^2$ is an isomorphism, so composing with the
unique corresponding deck shear makes it preserve $E_0$.  It then
preserves the torus $E_0\setminus\partial E_0$.  Evaluation at $(1,1)$
in the chart $(\mathbb C^*)^2\times\Delta_q$ shows that $a_i(q)$ are
holomorphic units at zero, and hence have logarithms.  Conversely their
multiplication extends on every toric chart and commutes with the deck
shears.  The logarithm kernel consists precisely of those periods with
no logarithmic term, namely $\ker(M_0-I)$.  On a finite root disc, a holomorphic lift has lattice defect in
$\ker N_j=(I-A_j)\Lambda$, by the saturated-image calculation
of \cref{fam:src:A.1.2}. Subtracting a period kills the defect
and preserves holomorphy. This gives the same exact sequence
with its finite equivariance.

\item   Put
\[
 M_0=(A_1A_2)^{-1}=\begin{pmatrix}1&0&0&0\\0&1&0&0\\0&1&1&0\\-1&0&0&1\end{pmatrix}.
\]
A cocycle on the punctured orbifold is a pair $(x_1,x_2)$ with
$N_1x_1=N_2x_2=0$, for the norm matrices in
\eqref{wb:pshift:eq:norms}.
Extension over the cusp says that $x_1+A_1x_2$ vanishes in
$\operatorname{coker}(M_0-I)$, so its first two coordinates vanish.
The complete integral solution is
\[
 x_1=(0,-b,2b-3c,c),\qquad x_2=(0,-2a-b,b,a),\quad a,b,c\in\mathbb Z.
\]
The principal cocycle of $(g,u,w,d)$ has parameters
\[
 \begin{pmatrix}a\\b\\c\end{pmatrix}
 =\begin{pmatrix}3&0&1\\-6&1&-1\\-2&1&0\end{pmatrix}
  \begin{pmatrix}g\\u\\w\end{pmatrix}.
\]
This matrix has determinant $-1$; its inverse is explicitly
$g=c-a-b$, $u=3c-2a-2b$, $w=4a+3b-3c$.  Hence
$H^1(\mathcal B,\mathcal H)=0$ integrally.

\item Sequence \eqref{eq:P5exponential} now makes every global translation a
global logarithm in $H^0(\mathcal B,E)=\mathbb C$.  Its kernel is
$\Lambda^{A_1,A_2}=\mathbb Z\widehat\delta$.  The flow extends by
commutation with the affine finite actions and by torus multiplication
on the cusp charts.  This proves existence, faithfulness, and
completeness of the asserted group.
\end{enumerate}
\end{proof}
